\documentclass[../main.tex]{subfiles}
\graphicspath{{\subfix{../images/}}}

\begin{document}

\newpage
\chapter{BACKGROUND}

This chapter mainly addresses the equations used by the one-dimensional finite-difference time-domain (FDTD) method, it's flow and the definitions of two differente technologies used to execute code in GPU's: Julia CUDA and Vulkan. Each one with a short description and a 1D FDTD implementation done with it.

\section{Yee's algorithm}

The finite-difference time-domain (FDTD) method\cite{ufdtd}, also known as Yee's algorithm, is used to aproximate electromagnetic fields over time.
It consists on a pair of ``update equations'' for the \blackcircleref{step:magnetic} magenetic fields and \blackcircleref{step:electric} electric fields based on discrete versions of Ampère's law and Faraday's law that are called $t$ times.
Each equation takes the previous iteration's fields and some constants from the material properties as inputs and sets the current iteration's magnetic or electric fields.

%The core of Fdtd Lucuma's simulation loop at each time step consists of the previously mentioned two equations, one or more additive electric sources and the boundary conditions of each field and component.
%Fdtd Lucuma currently only has the gaussian source as is sole source implementation and the Absorbing boundary condition as the boundary implementation.

\begin{figure}[H]
	\centering
	\begin{tikzpicture}[
			transform shape,
			node distance=1cm
		]

		\node[startstop] (start) {Start};

		\node[decision, below=of start, yshift=0.25cm] (while) {Can step?};

		\node[startstop, right=of while] (end) {End};

		\node[process, below=of while] (updateH) {\blackcircle{step:magnetic} Update magnetic $H_{\{x,y,z\}}$ fields};
		\node[process, below=of updateH] (updateE) {\blackcircle{step:electric} Update electric $E_{\{x,y,z\}}$ fields};
		\node[process, below=of updateE] (sources) {\blackcircle{step:sources} Apply sources};
		\node[process, below=of sources] (abc) {\blackcircle{step:abc} Apply boundary conditions};

		\draw[arrow] (start) -- (while);
		\draw[arrow] (while) -- node[anchor=south] {No} (end);
		\draw[arrow] (while) -- node[anchor=east] {Yes} (updateH);
		\draw[arrow] (updateH) -- (updateE);
		\draw[arrow] (updateE) -- (sources);
		\draw[arrow] (sources) -- (abc);
		\draw[arrow] (abc.west) -| ++(-2,0) |- (while.west);

	\end{tikzpicture}
	\caption{3D Simulation step flowchart}
	\label{fig:fdtd_lucuma_step}
\end{figure}

\subsection{1D equation}

For the sake of conciseness only the one-dimensional equations will be shown. The full 3D equations are in Appendix \ref{section:fdtd_3d}.

\subsection{Discrete notation}

\citet{ufdtd} used the following notation to represent the discretized versions of the electric $E$ and magnetic $H$ fields respectively.

\begin{equation}
	\overbrace{E_z(x,t)}^{\text{continuous}} = E_z(m\Delta_x, q\Delta_t) = \overbrace{E_z^q[m]}^{\text{discrete}}
\end{equation}
\begin{equation}
	H_y(x,t) = H_y(m\Delta_x, q\Delta_t) = H_y^q[m]
\end{equation}

Where $m,q$ are the spatial and temporal step indices, and $\Delta_x,\Delta_t$ are their respective offsets.
When using 3 spatial dimensions $m,n,p$ are the $x,y,z$ indices.

\subsection{Magnetic field update}

$\mu$ is the permeability of the material.

\begin{equation}
	\text{\blackcircleref{step:magnetic}\quad}
	H_y^{q+\frac{1}{2}} \left[m+\frac{1}{2}\right] = H_y^{q-\frac{1}{2}} \left[m+\frac{1}{2}\right] +\frac{\Delta_t}{\mu\Delta_x} \left(E_z^q\left[m+1\right]-E_z^q\left[m\right]\right)
\end{equation}

\subsection{Electric field update}

$\epsilon$ is the permittivity of the material.

\begin{equation}
	\text{\blackcircleref{step:electric}\quad}
	E_z^{q+1}\left[m\right] = E_z^{q}\left[m\right] + \frac{\Delta_t}{\epsilon\Delta_x}\left(H_y^{q+\frac{1}{2}}\left[m+\frac{1}{2}\right]-H_y^{q+\frac{1}{2}}\left[m-\frac{1}{2}\right]\right)
\end{equation}

\subsection{Gaussian source}

This source depends on the $q$ temporal step which is scaled by $\sigma$.
\begin{equation}
	\text{\blackcircleref{step:sources}\quad}
	E_z^{q+1}\left[m\right] = E_z^{q}\left[m\right] + \exp\left(-\left(\frac{q}{\sigma}\right)^2\right)
\end{equation}

\subsection{Absorbing boundary conditions}
$Cr$ is the Courant number.

\begin{equation}
	\text{\blackcircleref{step:abc}\quad}
	E_z^{q+1}\left[1\right] = E_z^{q}\left[1\right] + \frac{\frac{Cr}{\sqrt{\mu*\epsilon}}-1}{\frac{Cr}{\sqrt{\mu*\epsilon}}+1}*\left(E_z^q\left[2\right]-E_z^q\left[1\right]\right)
\end{equation}
\begin{equation}
	\text{\blackcircleref{step:abc}\quad}
	E_z^{q+1}\left[x\right] = E_z^{q}\left[x\right] + \frac{\frac{Cr}{\sqrt{\mu*\epsilon}}-1}{\frac{Cr}{\sqrt{\mu*\epsilon}}+1}*\left(E_z^q\left[x-1\right]-E_z^q\left[x\right]\right)
\end{equation}

\section{Julia + CUDA}
Julia\cite{juliaspec} is a dynamically typed programming language mainly used for scientific computing.
Moreover, the \texttt{CUDA.jl}\cite{besard2018juliagpu} package allows it to interface with Nvidia GPUs.
Here is an example of a 1D FDTD implementation.
\inputminted[bgcolor=codeBg, tabsize=4, breaklines=true, linenos=true]{julia}{\subfix{add.jl}}

\section{Vulkan}
Vulkan\cite{vulkanspec} is a low-level C API for graphics and compute in GPUs.
Unlike CUDA, which only runs on Nvidia hardware natively\footnote{With ZLUDA is posible to run CUDA programs in Radeon hardware.}, it has a wider range of compatible GPU hardware\cite{vulkandb}.

In OpenGL, Vulkan's predecessor, the shader source file can be passed directly to the driver, but as shown in Figure \ref{fig:vulkan_compute_pipeline} Vulkan requires
for each shader to be compiled into a spirv file, loaded into a shader module\footnote{From Vulkan 1.4 onward the \texttt{*.spv} file path can be passed directly into the creation of the pipeline.} and then finally used as part of a compute pipeline before it can be used.
For that purpose a shader compiler is needed.
Like Google's \texttt{glslc}\cite{shaderc} or Shader Slang's \texttt{slangc}\cite{10.1145/3197517.3201380} (which is the one that we used in this project).


\begin{figure}[H]
	\centering
	\begin{tikzpicture}[
			transform shape,
			vertex/.style={draw, ellipse, thick},
			scale=0.9
		]
		\node[vertex] (A) at (0,0) {\texttt{*.spv}};
		\node[vertex] (F) [left=of A] {\texttt{*.hlsl}};
		\node[vertex] (B) [above =of F] {\texttt{*.slang}};
		\node[vertex] (C) [below=of F] {\texttt{*.glsl}};
		\node[vertex] (D) [right=of A] {\texttt{VkShaderModule}};
		\node[vertex] (E) [right=of D] {\texttt{VkPipeline}};

		\path[-{Triangle[length=5]}, thick]
			(B) edge (A)
			(C) edge (A)
			(F) edge (A)
			(A) edge (D)
			(D) edge (E)
			(A) edge[bend left, dashed] node[label=above:Since Vulkan 1.4\cite{vulkan14}] {} (E)
		;

	\end{tikzpicture}
	\caption{Vulkan compute pipeline creation}
	\label{fig:vulkan_compute_pipeline}
\end{figure}

Using Vulkan Compute Shaders requires writing a considerable amount of boilerplate.
Therefore for the 1D example, only the \texttt{*.slang} file will be shown.

\inputminted[bgcolor=codeBg, tabsize=4, breaklines=true, linenos=true]{hlsl}{\subfix{updateH.slang}}
\inputminted[bgcolor=codeBg, tabsize=4, breaklines=true, linenos=true]{hlsl}{\subfix{updateE.slang}}
\inputminted[bgcolor=codeBg, tabsize=4, breaklines=true, linenos=true]{hlsl}{\subfix{gaussianSource.slang}}
\inputminted[bgcolor=codeBg, tabsize=4, breaklines=true, linenos=true]{hlsl}{\subfix{abc.slang}}

%\inputminted[bgcolor=codeBg, tabsize=4, breaklines=true, linenos=true]{cpp}{\subfix{add.cpp}}


\end{document}
